.bp73
.mt9
.mb7
.po8
.ll70
.ig
CHAPTER 27
D M S	  D I S P L A Y    F O R M A T S
.li
		     Your Company Name

	 I N V E N T O R Y   C O N T R O L   S Y S T E M

[  /  /  ]	      PROCESSING CYCLE

[ ]  1 ADD TO STOCK		     [ ]  4 REORDER REPORT
[ ]  2 DEPLETE STOCK		     [ ]  5 VALUATION REPORT
[ ]  3 TRANSACTION PROOF	     [ ]  6 ACTIVITY REPORT
				     [ ]  7 START ACTIVITY PERIOD

		       SYSTEM STARTUP

[ ]  8 PARAMETER ENTRY		     [ ]  9 CREATE SORT PARAMETER FILES

		   ITEM FILE MAINTENANCE

[ ] 10 ITEM ENTRY		     [ ] 12 ITEM FILE PRINT
[ ] 11 ITEM FILE SORT		     [ ] 13 ITEM FILE RECOPY

		   SELECT A PROGRAM > [  ]

***********************************************************************

		      Your Company Name
      00/00/00	  P A R A M E T E R   E N T R Y   M E N U

		  1  GENERAL PARAMETERS
		  2  AUDIT PARAMETERS

		  SELECT A SCREEN > [  ]
.BR
.he Appendix: Chapter 27
.BP
.LI
		     Your Company Name

		     PARAMETER ENTRY
    00/00/00	G E N E R A L	 P A R A M E T E R S

COMPANY NAME [						     ]

CURRENT PERIOD [       ]	       AVERAGE VALUATION  [   ]
     (WEEK, MONTH, QUARTER)

TO DATE PERIOD [       ]	       ARRAY SIZE	[     ]
     (MONTH, QUARTER, YEAR)

DATE FORM		  [ ]	       CONSOLE BELL USED   [  ]
     (1 MM/DD/YY  2 DD/MM/YY)

PAGE WIDTH	[   ]		       DEBUGGING	   [  ]

LINES PER PAGE	[   ]
.br
.bp
.sp3
.LI
			 Your Company Name

			  PARAMETER ENTRY
	 00/00/00      A U D I T   P A R A M E T E R S

		       AUDIT DRIVE  [ ]

		       ADDITIONS AUDIT USED [ ]
		       DEPLETIONS AUDIT USED [ ]



****************************************************************


			 Your Company Name
  DATE [  /  /	]      INVENTORY CONTROL SYSTEM
			 I T E M   E N T R Y

  ACCESS KEY					    ITEM NUMBER
[	    ]					    [	      ]

DESCRIPTION	  [				 ]

STOCKROOM QUANTITY [	   ]		VENDOR NUMBER [     ]

REORDER POINT	   [	   ]		COST / UNIT   [       ]

STD REORDER QTY    [	   ]		QTY ORDERED [	     ]

RETAIL PRICE	   [	       ]	DATE ORDERED  [        ]

INVENTORY VALUE    [		 ]	DATE DUE      [        ]

					BACK ORDERED  [       ]
.br
.bp
.li
			Your Company Name
		     INVENTORY CONTROL SYSTEM

DATE [	/  /  ]      ADDITIONS TO STOCK 	BATCH NUMBER [	    ]


ITEM NUMBER   [ 	 ]		 QUANTITY ORDERED  [	    ]

QUANTITY ADDED [       ]		 DATE ORDERED [        ]

UNIT COST      [	  ]		 DATE DUE     [        ]

STOCKROOM QUANTITY [	    ]		 QUANTITY BACKORDERED [        ]

	    DESCRIPTION [				 ]


************************************************************************


		       Your Company Name
		   INVENTORY CONTROL SYSTEM

DATE [	/  /  ]    DEPLETIONS FROM STOCK	   BATCH NUMBER  [     ]



		  ITEM NUMBER	 [	  ]

		  QUANTITY DEPLETED [	     ]

		  STOCKROOM QUANTITY [	      ]

		  DESCRIPTION [ 			]
.br
.he
.bp
.mt9
.mb7
.po8
.ll65
.ig
CHAPTER 28
D M S	S E T U P    F O R   N O N - S T A N D A R D   C R T ' S
.pp
To create a new CRT.DEF file, run the CRTDEF program, answering all the
requests that pertain to the CRT you are using.  When the program requests
the name of the CRT definition file to create, respond with CRT.  The
program will create a file named CRT.DEF.  This file should be placed on your
program disk (which should be on the currently logged drive).
All of the other DEF files can be omitted from your work disks or system
master (but not from the SSG distribution disk), as well as the CRTDEF.INT
program.
.he Appendix: Chapter 28
.pp
The utility program CRTDEF is used to define the physical characteristics of
a CRT display device (e.g., number of rows and columns, screen clear
sequence, cursor addressing sequence, etc.).  The program will prompt for
values for all CRT control sequences which are supported by the DMS system.
.pp
The program should only have to be executed once to define the CRT you are
using.	It does not have to be kept on the sytem disk after it is run.
.pp
The CRTDEF program does not use screen formatting and will work on any serial
CRT.  It will request data for almost every conceivable type of CRT control
procedure.  Not all of these procedures are need by the Inventory Control
System programs.  The only control sequences that are necessary to run are
the direct cursor addressing control characters and translation tables,
destructive backspace, and the clear screen character.	Any other requests
can be ignored by pressing RETURN.
.pp
Some control items will cause the system to execute faster because it can take
advantage of more sophisticated hardware features.
.pp
If you are modifying an existing CRT definition file, the CRTDEF program
will present a menu so that you may choose which section of the program to run.
Enter the number of your desired selection.  If you are creating a CRT
defintion file for the first time, the menu does not appear and all of the
requests that are not applicable can be answered to with a RETURN.
.pp
When the program issues a request  the current value of the control sequence
is displayed and the user is prompted for either a new value or acceptance
of the current value.	To accept the current value, just press the RETURN
key.
.pp
When the program starts up, it will request the name of the CRT physical
characteristics definition file.  For the Inventory Control System, always
enter CRT.  This is the name of the definition file that is created by
CRTDEF and the DEF file the Inventory Control System will look for.
.pp
CRTDEF will then request whether this run will create the intial version
of the file or edit an existing copy.  The first time you run this
program enter a C for create.  If you are modifying a DEF file that was
entered incorrectly, enter an M for modify.
.cp8
.sp2
THE CONTROL SEQUENCE REQUESTS:
.pp
Most of the CRT control sequences are entered in the form:
.sp
.ti8
<number of items>,<item-1>,<item-2>, . . .
.pp
The values of the items are the data characters which will be transmitted
to the CRT device to effect the speciafied control operation.  These
values are specified as decimal integers between 0-255.  To enter an
ESCAPE, for example, use the decimal equivalent of the ASCII value which
is 27 in decimal.  If a control sequence consisted of an ESCAPE followed by
a left parenthesis, you would enter:
.sp
.ti8
2,27,40
.sp
.ad
Where the 2 indicates that there are two characters that must be sent to the
CRT for this function and the 27 and 40 are the decimal representations
of an ESCAPE and a left parenthesis.
.pp
Please note that the CRTDEF program does not check the values for validity.
Valid values range from 0 to 255 but any integer value may be supplied.  The
value modulo 256 will be used.
.pp
The values specified are stored in character strings and are written
to the disk file as such.  CBASIC and CP/M impose restrictions on the
values which may be contained in these strings.  The restrictions are:
.sp
.ti8
CP/M:  CTRL-Z (26) is used as an end-of-file marker.
.sp
.ti8
CBASIC:  A RETURN (10) or a LINE-FEED (13) character will denote an end of
line sequence.
.pp
The way around the restrictions for most terminals which might happen to use
these characters for CRT control is to specify the value as <value>+128.
This sets the ASCII parity bit to ON.  For example, to send a decimal
10, send 138 instead.
.pp
The ASCII terminals will ignore the parity bit because the defined
ASCII set is 7 bits.  (Normally the CBIOS [operating system] will force
the parity bit to a fixed value before transmitting it to the terminal,
so the problem never really arises.
.sp2
REQUESTS ISSUED BY CRTDEF:
.sp
NUMBER OF ROWS FOR DISPLAY:
.in8
This is the number of rows (lines) of vertical display on your CRT.
.sp
.li
NUMBER OF COLUMNS FOR CRT DISPLAY:
.in8
This is the number of columns (characters across) displayed on your CRT.
.sp
.li
SPECIFY CURSOR ROW ADDRESS POSITIONING CONTROL CHARACTERS:
	row 1  current value:  0
		   new value:
.in8
Enter the value required by the CRT to position the cursor at the desired row.
Don't include the position control character sequence, merely the number
that will be translated into a row address.  This request will be issued
once for every possible row.
.sp
.li
SPECIFY CURSOR COLUMN ADDRESS POSITIONING CONTROL CHARACTERS:
	column 1   current value:  0
		       new value:
.in8
Enter the value required by the CRT to position the cursor at the desired
column.  Don't include the position control character sequence, merely the
number that will be translated into a column address.  This request will
be issued once for every possible column.
.sp
.li
CRT CONTROL:  SET CURSOR ADDRESS:
.in8
As described above, this request starts with "CRT control:" and thus requires
the number of elements to precede the values of the elements and for them to
be separated by commas.
.sp
Enter the values only if your CRT uses one contrl character sequence to send
the cursor to both row and column.  Enter the characters required to send the
cursor to a ROW and a COLUMN at the same time.
.bp
.li
SET ROW ADDRESS PREFIX:
.in8
If the CRT requires different control characters for ROW and COLUMN, enter
the row characters here (the previous request can be left a zero).  If the
CRT address row and column at the same time this can be left zero.
.sp
.li
SET COLUMN ADDRESS PREFIX:
.in8
If the CRT requires different control characters for ROW and COLUMN, enter
the column characters here.  If the CRT addresses row and column at the same
time this can be left zero.
.sp
.li
SPECIFY ORDER OF PRESENTATION OF ROW AND COLUMN ADDRESS:
.in8
Enter an "R" or a "C" for the order in which the control characters are to
be sent.
.sp
.li
SPECIFY FORMAT OF ROW AND COLUMN CURSOR CONTROL CHARACTERS:
.in8
Most CRT's require control information to be sent as hexadecimal characters,
but some require ASCII numbers.  Enter an "H" or a "D" for the correct format.
.sp
.li
CLEAR SCREEN:
.in8
Enter the characters for clearing the CRT screen.  This is a required entry.
.sp
.li
CLEAR FOREGROUND DATA FIELDS:
.in8
If your CRT supports different display intensities, enter the
foreground/background
control characters for this and the next 2 requests.
.sp
.li
ENTER FOREGROUND MODE:
.SP
.li
ENTER BACKGROUND MODE:
.SP
.li
RETURN CURSOR TO HOME POSITION:
.SP
.li
MOVE CURSOR UP ONE ROW:
.SP
.li
MOVE CURSOR DOWN ONE ROW:
.SP
.li
MOVE CURSOR LEFT ONE COLUMN:
.SP
.li
MOVE CURSOR RIGHT ONE COLUMN:
.SP
.li
DESTRUCTIVE BACKSPACE:
.in8
This sequence must be entered.	This is the sequence of characters
that will erase the previously typed in character.  This will
almost always be a backspace followed by a blank followed by a backspace.
In decimal, this translates to an 8, a 32, another 8.  Key in:
3,8,32,8.
.sp
.li
SOUND CONSOLE ALARM
.in8
This is normally a 1,7.
.sp
.li
ERASE TO END OF CURRENT LINE:
.SP
.li
ERASE TO END OF SCREEN:
.SP
.li
INSERT A NEW LINE INTO THE DISPLAY:
.SP
.li
DELETE CURRENT LINE FROM DISPLAY:
.SP
.li
NUMBER OF STORAGE LOCATIONS TO BE SET AT CRT INITIALIZATION:
.in8
Specify number of storage locations to be set or press RETURN.	If your
CRT requires memory to be initialized before operating, enter the number
of bytes that must be set at this request.
.sp
.li
CURRENT STORAGE ENTRY AND VALUES:   0	 0
.br
.li
SET STORAGE ENTRY 1:
.in8
Specify memory location and value.  This request will be issued for each location
specified in the previous request.  Enter the decimal value for each byte.
.sp
.li
INITIALIZATION SUBROUTINE ADDRESS:
.in8
If a machine code routine resides in memory that must be branched to on
startup, enter its address in decimal.
.sp
.li
INITIALIZATION I/O PORT:
.in8
If data needs to be transmitted to an I/O port, enter its number.
.sp
.li
DATA TO BE TRANSMITTED TO CONSOLE:
.in8
If data must be sent to the console device itself at startup time, enter
the number of bytes and their decimal values here.
.qi
.sp
.pp
The requests which follow, or
options 8 and 10 on the CRTDEF menu are provided to allow you to change
the normal definition of the special function control keys in the system.
If, for example, your operating system or computer does not allow free use
of a CONTROL-A character (enter adding mode), the CRTDEF program would
allow you to define a CONTROL-Y entered at the keyboard to be interpreed
by the program as though you had entered a CONTROL-A.  Because the
facilitiy is completely generalized, it should be possible to compensate
for the special requirements of all of the CP/M derivative operating
systems.
.pp
In addition, the use of control keys can be completely bypassed if
necessary.  A prefixing character can be defined that substitutes for
a control key.	In this manner, for example, an at-sign (@) or an
up arrow (^) can be designated as the prefix character.  Then an
@A sequence  would be interpreted the same as a CONTROL-A.
.pp
The option #8 requests
allow a control character translation table to be entered.  It
issues up to 31 requests for each control character possible.  If, for
example, your system has forced you to not use a CONTROL-A (ASCII value
of 1) and you will substitute a CONTROL-Y (ASCII value of 25), when the
request:
.sp
.li
	CHARACTER  25	    CURRENT TRANSLATION   25
	ENTER NEW VALUE (OR STOP) >
.br
.sp
.ad
is issued, respond with the value 1.  Don't forget that 8 is the backspace
and 13 is the carriage return.
.pp
The option #10 requests are
used only if you cannot use any control characters at all.
They will define instead a prefix character and the individual values
represented by the prefix character followed by one of the 127 possible
ASCII values.
.pp
The first request issued is for the value of the prefix character itself.
Enter its decimal ASCII value.	An at-sign, for example, would be 64.
Remember that once a character is defined as the prefixing character it
cannot be used as a valid data character anywhere else.
This is very important, so choose carefully.
.pp
The next request is used merely to quickly fill all 127 elements of the
table with a default.	The default should be some invalid character
combination such as a zero.
.pp
The program will then issue up to 127 requests for translation of individual
characters.  If you wish for the character sequence "prefix-A" (65). Don't
forget to enter the same interpred value at the lower-case value of "A" (97).
.pp
If you use the prefix character facility of option 10, you should use the
control character translation facility of option 8 to "mask" out all special
function control characters.  This involves setting them to some harmless
but invalid value, such as a zero or 10 (LINE FEED).  All programs will
then reject this character with an explicit error message.
.pp
After all requests have been answered (when the menu is not used) or after
each menu funtion (when the menu is used), the program will display:
.sp
.li
	Specify S(ave), E(dit), or Q(uit)
.pp
If all the data was entered correctly, enter an S.  If the data needs
to be edited, enter an E and the menu will appear, allowing the
data just entered to be modified.  If you are unsure about the entries
appropriate for your device, see your dealer.
.he
.bp
.ig
CHAPTER 29
E R R O R    M E S S A G E S
.sp2
.ce
CBASIC ERROR MESSAGES
.pp5
SSG software is designed to trap all errors which can arise under normal
circumstances, but due to language limitations and/or hardware problems,
some errors may occur which the programs themselves cannot handle.  These
errors are CBASIC errors.  They appear in the following form,
.he Appendix: Chapter 29
.sp
.ti10
ERROR aa
.sp
.ad
where aa is two alphabetic characters issued by the CBASIC runtime package
CRUN2.	The occurance of a CBASIC error message usually indicates a serious
problem in your INT or data files.
.in8
.sp
.ti-8
ERROR RE:  This is caused by a record in a file which has too few or too
many file delimiter characters (quotes or commas), or whose length is
incorrect.  When a program then attempts to read that record, what
is on the disk does not match the record description, and CRUN2 issues
this message.
.sp
It is possible for this problem to occur because of various hardware problems
such as a bad write, a voltage surge, etc., but far and away the most
common cause  is a bad RAM location, which either changes some character
into a quote or a comma, or changes one of these file delimiters into
some other character.  This error is virtually never the result of a
program bug.  Once a bad record is in the file, it must be repaired or
removed with the editor, or the entire file must be rebuilt.  Bad records
can be detected by TYPE'ing the file and manually searching for a missing
or extra quote or comma.  If you are not familiar with the use of the
editor, read the CP/M
documentation before attempting to change the file.  Also
read the File Definition in the Inventory Control System Specifications
Manual, because
some files may require that the header record be changed if a record
is manually deleted.
.sp
.ti-8
ERROR IV:  This error occurs if you try to run a program compiled by
CBASIC version 1 with CRUN2, or if the INT file that you are trying to run
is null.  If the latter is the case, recopy the file from the original disk,
using the verify option [vo] of PIP.
.sp
.ti-8
ERROR DZ:  This error occurs if you try to divide by zero, or if you try
to run a program compiled under CBASIC2 with CRUN version 1.  Since none
of the Inventory Control System programs will allow a division with a
denominator of zero, you should verify that your runtime package is indeed
CRUN2 (ideally version number 2.04 or higher).
.sp
.ti-8
ERROR OM or GARBAGE ON SCREEN DURING AN ENTRY PROGRAM:	This means that the
program has run out of memory--should not occur if you have at least 48K
of RAM and a standard-sized version of CP/M.
.qi
.sp3
.ce
INVENTORY CONTROL SYSTEM ERROR MESSAGES
.sp
.pp5
The terms "part" and "item" are used interchangeably throughout
the documentation.  Messages that instruct you to contact your distributor
generally indicate an error condition that should not arise under normal
circumstances.
.sp
.ce
SYSTEMWIDE MESSAGES
.sp
.in6
.ti-6
IY000 PROGRAM MUST BE RUN FROM MENU: To enter the Inventory Control System,
type CRUN2 IY.
.sp
.ti-6
IY002 INVALID DRIVE SPECIFIED:	Must be a single alphabetic character.
.sp
.ti-6
IY005 INVALID DATE:  MM/DD/YY or DD/MM/YY only.
.sp
.ti-6
IY044 INVALID COMMAND SEQUENCE: Contact your distributor.
.sp
.ti-6
IY045 INVALID COMMAND SEQUENCE: Contact your distributor.
.sp
.ti-6
IY086 NO VALID PART RECORDS: Build the item file through the ITEM ENTRY
menu function.
.sp2
.ce
ITEM ENTRY (PARTENT.INT)
.sp
.in6
.ti-6
IY101 PART FILE FOUND FROM ABORTED RUN:  A part file was found whose type had
a final digit of 0.  This occurs if a program which writes to a part file
ends abnormally.  Determine the reason, rename the
part file (or files) so that its last digit is 1, and continue. See Chapter
30.
.sp
.ti-6
IY102 PART FILE NOT FOUND:  Mount the correct disks on the appropriate drives
and continue.
.sp
.ti-6
IY103 PART FILE MOUNTED ON WRONG DRIVE:  See IY102.
.sp
.ti-6
IY104 UNEXPECTED END OF FILE ON PART FILE:  The part file is null or the
header record incorrect.  Use the TYPE command to check for the former.
The simplest recovery procedure is to bring your backup file up to date.
See Chapter 30.
.sp
.ti-6
IY105 not used
.sp
.ti-6
IY106 RENAME FAILED: A rename failed during processing.  There should be no
damage to the file.  See Chapter 30 for how to rename the file.
.sp
.ti-6
IY107 A QUOTE IS AN INVALID CHARACTER: No double quotes allowed in any file.
.sp
.ti-6
IY108 INVALID PART NUMBER: Must be A-Z, 0-9,
a slash (/), or a dash (-).
.sp
.ti-6
IY109 NOT AN INTEGER: 0-9 only, no decimals or commas.
.SP
.TI-6
IY110 NOT NUMERIC: Not a numeric character, or
too many digits to the left or to the right of the decimal.
See documentation for size limits.
.SP
.TI-6
IY111 THAT'S NOT A VALID RESPONSE: You entered an invalid control character
for that field. See documentation for appropriate response.
.SP
.TI-6
IY112 RENAME FAILED AT END: One or more item files will need to be renamed.
There should be no damage to the file.	See Chapter 30 for how to rename the
file.
.sp
.ti-6
IY113 INVALID PART FILE KEY: For the access key @D(rive)N(umber), D is
not a valid drive reference, or N a valid record number.
.sp
.ti-6
IY114 PART NUMBER NOT FOUND: Either the file is not in sorted order, or
the part specified is not on the file.
.sp
.ti-6
IY115 A SORT WILL BE REQUIRED IF THIS PART IS ADDED: Only the file affected
will need to be sorted.  See Chapter 22 for sorting procedures.
.sp
.ti-6
IY116 PART NUMBER ALREADY IN FILE: No duplicate part numbers allowed.
.sp
.ti-6
IY117 NO PR2 FILE FOUND: Your system disk is not in the drive where it was
at the start of the current run.  If not, contact your distributor.
.sp
.ti-6
IY118 RECORD NUMBER OUT OF RANGE: For the access key @D(rive)N(umber),
N is not a valid record number on drive D.
.sp
.ti-6
IY119 INVALID PART NUMBER: see IY108
.sp
.ti-6
IY120 NO PR2 FILE FOUND: Contact your distributor.
.sp
.ti-6
IY121 THERE ARE NO VALID PART RECORDS. YOU'LL HAVE TO ADD SOME: Should
appear only at system startup.
.sp
.ti-6
IY122 THE PART RECORD IS BEING DELETED: informational message.
.sp2
.ce
PARAMETER ENTRY (IYPARENT.INT)
.sp
.ti-6
IY201 THIS PROGRAM MUST BE RUN FROM THE IY MENU: Type CRUN2 IY to enter the
Inventory Control System.
.sp
.ti-6
IY202 INVALID INPUT: Invalid control characters or an invalid response.
Consult the documentation for the appropriate response.
.SP
.TI-6
IY203 not used
.SP
.TI-6
IY204 INVALID DRIVE SPECIFIED: System drive must be either A or B.
.SP
.TI-6
IY205 not used
.SP
.TI-6
IY206 NO PARAMETER FILE. WILL CREATE DEFAULT FILE BEFORE CONTINUING:
This message should only appear at system startup.  If you do receive it,
stop the program and locate the parameter file, or allow the program to
create one.
.SP
.TI-6
IY207 UNEXPECTED END OF PARAMETER FILE: Parameter file not found when stop
requested. Contact your distributor.
.SP2
.CE
ADD TO STOCK (ADDITION.INT)
.SP
.TI-6
IY301 PART FILE FOUND FROM ABORTED RUN: see IY101
.sp
.ti-6
IY302 PART FILE NOT FOUND: see IY102
.sp
.ti-6
IY303 PART FILE MOUNTED ON WRONG DRIVE: see IY102
.sp
.ti-6
IY304 UNEXPECTED END OF FILE ON PART FILE: see IY104
.sp
.ti-6
IY305 not used
.sp
.ti-6
IY306 not used
.sp
.ti-6
IY307 not used
.sp
.ti-6
IY308 PR2 FILE NOT FOUND: The program was unable to find the PR2 file, or
found a null one.  Either PIP a good PR2 file onto the system disk, or
recreate one with the Parameter Entry program. This should only if disks are
switched improperly during ITEM FILE RECOPY or sorting.
.sp
.ti-6
IY309 TOO MANY DIGITS: Your entry was not numeric, or contained too many
digits.  The fields for quantity added, quantity on order, or quantity on
back order accept up to 6 digit entries.
.SP
.TI-6
IY310 INVALID RESPONSE: An inappropriate control character was entered.
See the documentation for the valid response.
.SP
.TI-6
IY311 INVALID RESPONSE: re-enter appropriate data.
.sp
.ti-6
IY312 QUANTITY IN STOCK MAY NOT BE LESS THAN ZERO
.SP
.TI-6
IY313 CURRENT PERIOD STOCK ADDITIONS CANNOT GO BELOW ZERO
.SP
.TI-6
IY314 CURRENT PERIOD STOCK ADDITIONS CANNOT EXCEED 999999
.SP
.TI-6
IY315 TO DATE STOCK ADDITIONS CANNOT GO BELOW ZERO
.SP
.TI-6
IY316 TO DATE STOCK ADDITIONS CANNOT EXCEED 999999
.SP
.TI-6
IY317 CANNOT BE LESS THAN ZERO: The quantity on order or back-ordered cannot
be less than zero or greater than 999999.
.SP
.TI-6
IY318 CANNOT EXCEED 999999: An attempt was made to either reduce an activity
accumulator below zero or increase one beyond its maximum.  The
accumulator will be set to zero or its maximum by the program.
.sp
.ti-6
IY319 INVALID DATE: Dates must be of the form MM/DD/YY or DD/MM/YY,
depending on what was specified in the Parameter Entry program.
.sp
.ti-6
IY320 AUDIT FILE CREATE FAILED: The ADD TO STOCK program
was unable to create the audit
file.  Perhaps the disk is full, or is write-protected.
.sp
.ti-6
IY321 AUDIT FILE FOUND FROM PREVIOUS ABORTED RUN:  The ADD TO STOCK or
DEPLETE STOCK program ended abnormally on its previous run, leaving an audit
file of type 000.  Find and correct the cause of the aborted run.  If the
audit file appears to be complete, rename it to type 001 and continue,
otherwise, erase it.  See Chapter 30.
.sp
.ti-6
IY322 AUDIT FILE OPEN FAILED:  The ADD TO STOCK program
was unable to open an audit file
which was listed in the directory, or found a null file.  Erase the file
and continue. See Chapter 30.
.sp
.ti-6
IY323 Y, N, OR ESCAPE ONLY: Re-enter appropriate data.
.sp
.ti-6
IY324 NOT NUMERIC: Re-enter appropriate data.
.sp
.ti-6
IY325 NULL AUDIT FILE FOUND: Possible hardware problem.
.sp
.ti-6
IY326 PART VALUE SET TO 9,999,999.99: This is the maximum inventory value
allowed for an item.  If it is not large enough for your needs, congratulations
on your good fortune.
.sp
.ti-6
IY327 AUDIT FILE OPEN FAILED AFTER SAVING FILES: see IY322
.sp
.ti-6
IY328 INVALID PART NUMBER: The part number specified does not exist in
the part file(s). Re-enter a valid part number: A-Z, 0-9, slashes (/), or
dashes (-) only.
.sp
.ti-6
IY329 INVALID PART KEY: A part key must be of the form @Dnnn, where
D is the drive reference for the desired part file, and nnn is the record
number of the desired part within that file.
.sp
.ti-6
IY330 VALUE SET TO ZERO: The transaction would take the value below zero.
.sp2
.CE
DEPLETE STOCK (DEPLETE.INT)
.sp
.ti-6
IY351 PART FILE FOUND FROM ABORTED RUN: see IY101
.sp
.ti-6
IY352 PART FILE NOT FOUND: see IY102
.sp
.ti-6
IY353 PART FILE MOUNTED ON WRONG DRIVE: see IY102
.sp
.ti-6
IY354 UNEXPECTED END OF FILE ON PART FILE: see IY104
.sp
.ti-6
IY355 INVALID DEPLETION QUANTITY: Your entry was not numeric, or its length
was greater than 6.
.sp
.ti-6
IY356 PR2 FILE NOT FOUND: see IY308
.SP
.TI-6
IY357 not used
.SP
.TI-6
IY358 PR2 FILE NOT FOUND: see IY308
.SP
.TI-6
IY359 UNEXPECTED EOF ON AUDIT FILE: Contact your distributor.
.sp
.TI-6
IY360 INVALID RESPONSE: Valid entries include CTRL-N, CTRL-B, ESCAPE,
and, if a valid
record has been read, RETURN, CTRL-R, and CTRL-S.
.SP
.TI-6
IY361 PART VALUE SET TO MIN. (ZERO): The depletion based on the current
average value per unit has forced the inventory value to zero.
.SP
.TI-6
IY362 STOCK MAY NOT BE LESS THAN ZERO: Your entry would set the quantity
on hand below zero.
.SP
.TI-6
IY363 CURRENT DEPLETIONS AT MIN. (ZERO): informational message.
.SP
.TI-6
IY364 CURRENT DEPLETIONS AT MAX. (9,999,999): informational message.
.SP
.TI-6
IY365 TO DATE DEPLETIONS AT MIN. (ZERO): informational message.
.SP
.TI-6
IY366 TO DATE DEPLETIONS AT MAX. (9,999,999): informational message.
.SP
.TI-6
IY367 PART VALUE SET TO MAX. (9,999,999): When backing out a depletion, the
inventory value has exceeded the maximum, based on the current average value.
.SP
.TI-6
IY368 STOCK MAY NOT EXCEED 999,999: Your entry would cause the stock on hand
quantity to exceed the maximum.
.SP
.TI-6
IY369 INVALID RESPONSE: This field will accept only a valid depletion quantity,
CTRL-B, CTRL-R, CTRL-X, or RETURN.
.SP
.TI-6
IY370 AUDIT FILE CREATE FAILED: Your disk or directory may be full.  If not,
contact your distributor.
.SP
.TI-6
IY371 AUDIT FILE FOUND FROM PREVIOUS ABORTED RUN: An IYAUD.000 file has been
found, meaning the ADD TO STOCK  or DEPLETE STOCK program ended abnormally.
Determine the reason for the abnormal ending, rename the file to type 001,
and continue.  See Chapter 30.
.SP
.TI-6
IY372 AUDIT FILE OPEN FAILED: Contact your distributor.
.SP
.TI-6
IY373 Y, N, OR ESCAPE ONLY: Y, N, T, True, Ok, or ESCAPE only.
.SP
.TI-6
IY374 not used
.SP
.TI-6
IY375 not used
.SP
.TI-6
IY376 not used
.SP
.TI-6
IY377 OPEN FAILED: Contact your distributor.
.SP
.TI-6
IY378 INVALID PART NUMBER: See IY328.
.SP
.TI-6
IY379 INVALID PART KEY: See IY329.
.SP2
.CE
START AN ACTIVITY PERIOD (UPDATE.INT)
.SP
.TI-6
IY401 PART FILE FOUND FROM PREVIOUS ABORTED RUN: see IY101
.sp
.ti-6
IY402 PART FILE NOT FOUND: see IY102
.SP
.TI-6
IY403 PART FILE MOUNTED ON WRONG DRIVE: see IY102
.sp
.ti-6
IY404 UNEXPECTED END OF FILE ON PART FILE: see IY104
.sp
.ti-6
IY405 PR2 FILE NOT FOUND OR NULL: see IY308
.SP
.TI-6
IY406 ACTIVITY REPORT HAS NOT BEEN RUN: Run the Activity Report, if desired,
and continue.
.sp
.ti-6
IY407 INVALID RESPONSE: Inappropriate entry.  See the documentation for
the proper response.
.SP
.TI-6
IY408 Y, YES, N, OR NO ONLY: Re-enter appropriate data.
.sp
.ti-6
IY409 not used
.sp2
.CE
IY MENU (IY.INT)
.SP
.TI-6
IY421 NO PARAMETER FILE: Contact your distributor.
.SP
.TI-6
IY422 DMS ERROR: CRT.DEF is an invalid definition file.  Contact your
distributor.
.SP
.TI-6
IY423 ERROR DURING CHAIN TO SORT: Disk or directory full, or disk is write
protected.  See Chapter 30.  If none of these, contact your distributor.
.SP
.TI-6
IY424 ERROR DURING CHAIN TO SORT: see IY423
.SP
.TI-6
IY425 THE PROGRAM IS NOT ON THE CURRENTLY LOGGED DISK: All INT files should
be located on the currently logged drive.
.SP
.TI-6
IY426 THE SORT PARAMETER FILE IS NOT ON THE SYSTEM DISK: Run the CREATE SORT
PARAMETER FILES  function.
.SP
.TI-6
IY427 ERROR DURING CHAIN TO SORT: see IY423
.SP
.TI-6
IY428 not used
.SP
.TI-6
IY429 THIS PROGRAM CAN ONLY BE RUN ON VERIFIED SORTED PART FILES:  All item
files must be sorted and checked for duplicate records (by running the ITEM
FILE PRINT program) before the program
selected can be run.
.SP
.TI-6
IY430 INVALID SECOND PARAMETER FILE: Invalid information in the IYPR2
parameter file. See Chapter 30.
.SP
.TI-6
IY431 QSORT IS NOT ON THE CURRENTLY LOGGED DISK: QSORT	is expected on the
currently logged disk.	See documentation
Chapter 22 for sorting procedures.
.SP
.TI-6
IY432 THAT'S NOT A VALID RESPONSE: The menu selection must be numeric.
ESC or CTRL-R (Refresh the screen) are the only other valid commands.
.SP
.TI-6
IY433 NO SECOND PARAMETER FILE. CHAINING TO PARTFILE: Since no parameter
file exists, the program
will create one and allow you to change the default values.
PARTFILE is the program responsible for the ITEM FILE RECOPY function.
.SP
.TI-6
IY434 THE SYSTEM DRIVE HAS NOT BEEN INITIALIZED: You cannot continue until
you specify either drive A or drive B as the system drive (the drive where
the parameter and SRT files are expected).
.SP
.TI-6
IY435 INVALID DATE: DD/MM/YY or MM/DD/YY only.
.sp
.ti-6
IY436 INVALID PARAMETER FILE: IYPR1 contains invalid information. See
Chapter 30.
.SP
.TI-6
IY437 OPERATOR REQUESTED RETURN TO MENU: informational message.
.SP
.TI-6
IY438 **ERROR**  THE PROGRAM WAS TERMINATED BEFORE COMPLETION: A fatal error
occured in the program just run. Determine the cause,
if possible, and continue.  See Chapter 30.
.SP
.TI-6
IY439 TRANSACTION HISTORY IS NOT CURRENTLY RECORDED: Auditability was not
requested.  Run PARAMETER ENTRY program to select it if desired.
.SP
.TI-6
IY440 not used
.sp2
.ce
REORDER REPORT (REORDER.INT)
.sp
.ti-6
IY441 PART FILE MOUNTED ON WRONG DRIVE: see IY103
.SP
.TI-6
IY442 UNEXPECTED EOF ON PART FILE: The part file was not found or contains
invalid information.  Find the partfile, or bring your backup file up to date.
.SP
.TI-6
IY443 THAT'S NOT A VALID RESPONSE: Only numeric entries allowed, no control
characters.
.SP
.TI-6
IY444 INVALID PART NUMBER: see IY108
.SP2
.CE
ITEM FILE PRINT (PARTPRNT.INT)
.SP
.TI-6
IY461 PART FILE FOUND FROM PREVIOUS ABORTED RUN: see IY101
.SP
.TI-6
IY462 PART FILE NOT FOUND: see IY102
.SP
.TI-6
IY463 PART FILE MOUNTED ON WRONG DRIVE: see IY102
.SP
.TI-6
IY464 UNEXPECTED END OF FILE ON PART FILE: see IY104
.SP
.TI-6
IY465 PART FILE RENAME FAILED: The program tried to rename an item file of
type 111 to 101 at end.
.sp
.ti-6
IY466 DUPLICATE PART NUMBER FOUND ON DRIVE: RUN PARTENT TO DELETE: Programs
in the Inventory Control System will not operate properly if there are
duplicate part numbers.  Delete the flagged duplicates with the ITEM ENTRY
program, and
re-run the ITEM FILE PRINT program.
.sp
.ti-6
IY467 not used
.SP
.TI-6
IY468 PR2 FILE NOT FOUND: see IY308
.sp
.ti-6
IY469 INVALID RESPONSE: Re-enter appropriate data.
.sp
.ti-6
IY470 INVALID PART NUMBER: A-Z, 0-9, slash (/), or dash (-) only.
.SP
.TI-6
IY471 TYPE 111 PART FILE FOUND WHEN TYPE 101 PRESENT: Remove 101 file.
.sp2
.CE
CREATE SORT PARAMETER FILES (SORTPARM.INT)
.SP
.TI-6
IY491 INVALID RESPONSE: Inappropriate entry.   See documentation for proper
response.
.SP
.TI-6
IY492 DRIVE REFERENCE MUST BE A-Z: Re-enter appropriate data.
.sp
.ti-6
IY493 UNABLE TO CREATE: Directory may be full or disk may not be write
enabled.
.sp
.ti-6
IY494 UNABLE TO WRITE: The CREATE SORT PARAMETER FILES program
was unable to create or write to the
specified file.  Perhaps the disk is full or write-protected.
.sp2
.CE
VALUATION REPORT (VALUE.INT)
.sp
.ti-6
IY501 INVALID PART NUMBER: see IY108
.sp
.ti-6
IY502 STARTING PART NUMBER GREATER THAN ENDING PART NUMBER: Re-enter
appropriate data.
.sp
.ti-6
IY503 PART FILE MOUNTED ON WRONG DRIVE: see IY102
.SP
.TI-6
IY504 UNEXPECTED END OF FILE ON PART FILE: see IY104
.SP
.TI-6
IY505 not used
.sp
.ti-6
IY506 INVALID RESPONSE: Re-enter appropriate data.
.SP2
.CE
ACTIVITY REPORT (ACTIVITY.INT)
.sp
.ti-6
IY511 PART FILE ON WRONG DRIVE: See IY102.
.sp
.ti-6
IY512 UNEXPECTED EOF ON PART FILE: See IY104.
.sp
.ti-6
IY513 PROGRAM MUST BE RUN FROM MENU: See IY201.
.SP
.TI-6
IY514 OPTION NUMBER OUT OF RANGE: Your entry was not numeric, greater than
6, or less than 1.
.SP
.TI-6
IY515 INVALID PART NUMBER: A-Z, 0-9, slash (/), or dash (-) only. A RETURN
defaults to the first part number on the file.
.SP
.TI-6
IY516 INVALID PART NUMBER: A-Z, 0-9, slash (/), or dash (-) only. A RETURN
defaults to the end of the file.
.SP
.TI-6
IY517 ENDING PART NUMBER IS LESS THAN STARTING PART NUMBER
.SP
.TI-6
IY518 PR2 FILE NOT FOUND: See IY308.
.SP2
.CE
TRANSACTION PROOF (PROOF.INT)
.SP
.ti-6
IY521 PROGRAM MUST BE CALLED FROM MENU: See IY201.
.SP
.TI-6
IY522 PROGRAM CAN'T RUN. NO AUDITING REQUESTED: Change the audit parameters
through the PARAMETER ENTRY program if auditability desired.
.SP2
.CE
ITEM FILE RECOPY (PARTFILE)
.SP
.TI-6
IY531 UNEXPECTED EOF ON PART FILE READ: Invalid information in part file
header record. See Chapter 30.
.SP
.TI-6
IY532 INVALID RESPONSE: Not a valid response.  See documentation for valid
responses.
.SP
.TI-6
IY533 INVALID PART FILE PARAMETERS: See documentation for proper file
allocation.
.SP
.TI-6
IY534 END OF FILE ON PART FILE CREATE: See Chapter 30.
.SP
.TI-6
IY535 THAT'S NOT THE RIGHT INPUT FILE: Put the correct file on the input
drive.
.SP
.TI-6
IY536 NO INPUT FILE FOUND ON DRIVE A: Put the correct item file on drive A.
.SP
.TI-6
IY537 INVALID PART FILE: The input file is null.  See Chapter 30.
.SP
.TI-6
IY538 THAT'S NOT THE RIGHT OUTPUT FILE: You've changed the B disk before
being prompted to do so.
.SP
.TI-6
IY539 UNEXPECTED EOF ON WRITE TO PART FILE: See IY423 and Chapter 30.
.SP
.TI-6
IY540 NO OUTPUT FILE FOUND: You've removed the B disk before being prompted.
.SP
.TI-6
IY541 INVALID OUTPUT FILE FOUND: The output file is null.  See Chapter 30.
.SP
.TI-6
IY542 THERE IS A PART FILE ON THE DISK ON DRIVE B: No part files should reside
on the disk used to receive the recopied output.  Mount a disk that does not
contain an IYPRT file.
.SP
.TI-6
IY543 THERE IS NO SECOND PARAMETER FILE ON THAT DISK: Check to make sure
your system disk (the disk with IYPR1 and IYPR2)
is in the drive named as the system drive (either A or B).
.SP
.TI-6
IY544 UNABLE TO CREATE PART FILE. IS THE DISK WRITE PROTECTED?: Same as IY423.
See Chapter 30.
.SP
.TI-6
IY545 UNABLE TO WRITE TO PART FILE: see IY544.
.SP
.TI-6
IY546 INVALID DATE: MM/DD/YY or DD/MM/YY only.
.SP
.TI-6
IY547 UNABLE TO CREATE PR2 FILE: see IY544.
.SP
.TI-6
IY548 NO INPUT FILE FOUND ON DRIVE A: Put item file to be recopied on drive A
as directed.
.SP
.TI-6
IY549 INVALID PART FILE: The file on the input drive is null.  See Chapter 30.
.SP
.TI-6
IY550 INVALID PART FILE: Files must be in sorted order.  If copying continues,
item files will not be in sorted order.
.qi
.he
.bp
.ig
CHAPTER 30
E R R O R    R E C O V E R Y
.pp5
When certain programs, notably the data and parameter entry, recopy, and
update programs end abnormally they may leave a file with an improper
filetype.  Under normal conditions the filetypes of item, parameter, and
audit files should end with a 1 (one) to indicate the file is a valid
current file (for example, IYAUD.001).	Since the 1 is changed to a 0 (zero)
when the file is being processed, and renamed back to 1 when processing
is complete, you will never see a zero as the third digit unless the program
ended before the rename took place, or the rename failed for some other reason.
.he Appendix: Chapter 30
.pp
Before renaming the file with the CP/M REN command, you should try to
learn why the program ended prematurely or the rename failed.  It is quite
possible that the problem is hardware related, and it may even be due to
a problem with the operating system.
.pp
To rename a file, give the REN command in the following form:
.sp
.in10
REN D:nnnnnnnn.ttt=D:nnnnnnnn.ttt
.qi
.sp
.ad
where D is a drive reference, nnnnnnnn an 8 character filename, and ttt the
3 character filetype.  The filename and type to the left of the equal sign
is that to which the file on the right will be renamed.  For example, to rename
IYPRT.100 to IYPRT.101 give the command,
.sp
.in10
REN A:IYPRT.101=IYPRT.100
.qi
.sp
.ad
if IYPRT.100 is on A, and you want IYPRT.101 on the A drive as well.  The
rename cannot take place if a file with the same name as that to which the
other file is to be renamed already exists on the designated drive.
.pp
Certain error conditions may arise because a disk or its directory is full,
or the disk write-protected.  If the disk is write protected, remove the
write-protect tab or apply a write-enable tab (whichever is appropriate for
your system).  Unless your disk contains many other files besides those used
and created by the Inventory Control System, the disk directory should not
be full.  However, since the number of files a directory can hold does
vary, you should find out the capacity of the directory you are using.
.pp
To check to see if a disk is full use the CP/M STAT command.  Although this
command should be on the currently logged drive, it can give you the status
of files on other drives.   For example, to find the number of bytes
remaining on the B drive with STAT on the A drive type the command,
.sp
.in10
STAT B:
.qi
.sp
.ad
See the CP/M Features and Facilities manual for complete instructions on
how to use the STAT command.
.pp
Some error messages make reference to a null file or null header record.
An INVALID FILE error message may also indicate a null file since no SSG
files are created as null files.  All SSG files have one record written
to them upon creation, either a header record, or if a parameter file,
the entire file is made up of one record.
If the entire file is null, the TYPE command, instead of displaying the
file contents, will simply return the CP/M prompt character.  If the header
record is null, the first record to appear will be a data record.  Currently
only the item and audit files use header records.
.pp
If a header record is incorrect, you may be able to correct it using a file
editor.  Check the file specifications for the necessary information.
Be certain to get the record length correct. This
recovery procedure is suggested only for those with substantial operating
experience or data processing knowledge.  In any case, back up your files
before attempting to modify them with an editor.
For others, the best solution is to
bring your backup file up to date.
.pp
To recover from certain error conditions you may need to recreate the first of
the two parameter files IYPR1.	Since no information in this file is
altered by processing, once the file has been created by the PARAMETER ENTRY
program you can continue.
.pp
If a recovery procedure requires you to recreate the second parameter file
IYPR2 you may need to alter the file to get back to where you were.  This
is because the IYPR2 file contains information that is altered during the
normal processing cycle.  You may want to reset the batch number to
correspond with the current number.  To do so, check the file specifications
and make the changes with an editor. The IYPR2 file also contains a
flag which indicates whether any files are unsorted.  You need not reset this
flag if you make sure all your item files are in sorted order before
continuing.
.bp
.pp
Message IY438 **ERROR** THE PROGRAM WAS TERMINATED BEFORE COMPLETION appears
at the system menu after a fatal error.  The error that caused IY438 to
appear should have displayed its error message briefly before the menu and
IY438 appeared.  If for some reason you did not notice which error it was,
restart the system and try to continue where you left off.  A rename may be
required if the program just run was either ADD TO STOCK, DEPLETE STOCK,
ITEM ENTRY, OR START A NEW ACTIVITY PERIOD.
If the error
occurs again, be sure to take note of the message.
.he
.bp
.ig
CHAPTER 31
C O M M A N D	 S U M M A R Y
.in9
.ti-9
ESCAPE:  Ends the current function or program.
.sp
.ti-9
RETURN:  Enters data and advances to next request, accepts default value,
moves cursor forward through DMS screen fields.
.sp
.ti-9
STOP:	 Typing S-T-O-P followed by RETURN is the same as ESCAPE, use at
non-DMS screen requests only.
.sp
.ti-9
CTRL-H:  Backs cursor up one position, erasing the character in that position.
.sp
.ti-9
DELETE (SHIFT-DELETE):	Backs up cursor one position, erasing the character
in that position; use at non-DMS requests.
.sp
.ti-9
CTRL-R:  Refreshes garbled DMS screen.
.sp
.ti-9
CTRL-X:  Cancels current add to stock, deplete stock, or item entry
transaction.
.sp
.ti-9
CTRL-A:  At ITEM ENTRY access key field, puts program in add-records mode.
.sp
.ti-9
CTRL-D:  At ITEM ENTRY access key field, deletes record currently displayed.
.sp
.ti-9
CTRL-S:  At ITEM ENTRY access key field, saves records added or changed so
far this session.
.sp
.ti-9
CTRL-N:  At ITEM ENTRY access key field, if no record currently displayed,
gets first record on file; if a record is currently displayed, gets next
record after it.
.sp
.ti-9
CTRL-B:  At ITEM ENTRY access key field, gets record before currently
displayed record; within a DMS screen, moves cursor back one field.
.sp
.ti-9
^ (UP ARROW):  Moves cursor back one request at non-DMS screen request.
.qi
#
